\section{Related Work and Research Gaps}
\label{sec:related}

We searched Crossref, arXiv, Europe PMC, PubMed Central, OpenAlex and publisher pages in October 2026, using keywords (waste reviews, datasets and detectors; India and South Asia; metaheuristic and hyperparameter optimization) and following citations from six seed papers~\cite{majchrowska2022deep,yang2016classification,kaya2023classification,sayed2024intelligent,aguerchi2024cnn,alshammari2025smart}. We kept vision papers on waste or litter and optimization papers on CNNs, mostly peer-reviewed work from 2016 to 2026. The review is targeted and not systematic, so statements about ``the literature'' below refer to the papers we reviewed. We read 15 model or dataset studies in full (Table~\ref{tab:lit-matrix}). For some of the other works we had only the abstract, and we cite those only for what the abstract states.

\subsection{Vision-Based Waste Recognition}
\label{subsec:rw-waste}

Much of the work on waste sorting classifies images of single objects. TrashNet, released with a Stanford course report by Yang and Thung~\cite{yang2016classification}, has 2,527 images in six classes, each showing one item on a white posterboard. In that report a support vector machine on SIFT features reached 63\% test accuracy and a CNN trained from scratch only 22\%. The authors put this down to the hyperparameters and left a proper hyperparameter search and multi-object detection for future work. Kaya~\cite{kaya2023classification} truncates pretrained Xception and InceptionResNetV2 networks and reaches 89.72\% test accuracy with one fixed training setting, and names TrashNet's plain backgrounds as a limitation. Sayed et al.~\cite{sayed2024intelligent} tune InceptionV3 with a metaheuristic (Section~\ref{subsec:rw-hpo}). The smart bin of Alshammari et al.~\cite{alshammari2025smart} sorts items with a MobileNetV3-based classifier on a Raspberry Pi, and TrashBox~\cite{kumsetty2022trashbox} is a classification set of 17,785 images in seven classes, although its title says ``detection''. One label per image is enough for a chute that takes one item at a time. It is not enough for a heap, a conveyor or a street.

The review papers say much the same. Abdallah et al.~\cite{abdallah2020artificial} survey 85 studies of artificial intelligence in solid waste management from 2004 to 2019. Lu and Chen~\cite{lu2022computer} describe many waste-sorting vision studies as preliminary, set in simplified environments and built on artificially collected data, and ask for more realistic scenes and open datasets. Fang et al.~\cite{fang2023artificial} find that most vision studies they cover classify whole images by material, that detection studies are fewer and less accurate, and that lab pictures differ from real waste. The accuracies they collect (72.8\% to 99.95\%) mix classification and detection, so they cannot be used to rank methods.

Detection and segmentation datasets do exist for cluttered scenes. TACO~\cite{proenca2020taco}, a crowd-sourced dataset published as a preprint, had 1,500 photographs of litter in the wild with 4,784 annotations at release. ZeroWaste~\cite{bashkirova2022zerowaste} shows deformable and translucent objects in heavy clutter on a sorting line, and WaRP~\cite{yudin2024hierarchical} targets hierarchical detection in recycling-plant images. UAVVaste~\cite{kraft2021autonomous} has 772 low-altitude drone images with 3,716 annotations; its authors found TACO's large, near-ground objects a poor fit for aerial detection. PlastOPol~\cite{cordova2022litter} has 2,418 urban, beach and forest images with 5,300 annotations of a single ``litter'' class. The TrashCan report~\cite{hong2020trashcan} labels 7,212 deep-sea images with boxes and masks, FloW~\cite{cheng2021flow} covers floating waste in inland waters, and MJU-Waste~\cite{wang2020multilevel} has 2,475 RGB-D lab images of people holding waste, again with one ``waste'' class. Each setting (aerial, underwater, conveyor, hand-held, street) has its own data, often with one generic class, and the class lists do not match across datasets.

The pipeline closest to ours is that of Majchrowska et al.~\cite{majchrowska2022deep}. They merge public datasets into a one-class ``detect-waste'' set and a seven-category ``classify-waste'' set that follows the sorting rules of Gda\'{n}sk. Direct multi-class detection worked poorly for them, so they localize litter with EfficientDet-D2 and classify each crop with EfficientNet-B2. The hyperparameters were set by hand, and results are single runs on a held-out 20\% that they describe as a validation and testing set. Other detection studies include Fulton et al.~\cite{fulton2019robotic}, who compare YOLOv2, Tiny-YOLO, Faster R-CNN and SSD for underwater litter; Melinte et al.~\cite{melinte2020deep}, who add boxes to TrashNet and report 97.63\% precision with SSD; and Liao and Juang~\cite{liao2022realtime}, who run YOLOv4-Tiny-3l at 22 frames per second on an Xavier NX for coastal litter monitoring by drone. Bartolo et al.~\cite{bartolo2025learning} train five detectors on one drone dataset and test them on two others, and S\'{a}nchez-Ferrer et al.~\cite{sanchezferrer2023experimental} study Mask R-CNN for marine debris when data are scarce, adding synthetic scenes. In all of these the hyperparameters are mentioned only as part of the setup (Table~\ref{tab:lit-matrix}).

Our detector belongs to the single-stage family that began with YOLO~\cite{redmon2016you}, which is the usual choice in waste studies that care about speed~\cite{kumar2021novel,kraft2021autonomous,liao2022realtime,cordova2022litter}. In the PlastOPol study, YOLOv5x was the most accurate of seven detectors and YOLOv5s ran at up to 5.19 frames per second on a phone. We use YOLOv8n~\cite{jocher2023ultralytics} with COCO-pretrained weights~\cite{lin2014microsoft} (Section~\ref{subsec:detector}).

\subsection{Waste Data From India and South Asia}
\label{subsec:rw-india}

Waste image data from India are scattered, and most were not built for detection. The earliest we found is SpotGarbage~\cite{mittal2016spotgarbage} from IIT Ropar, a smartphone app that detects and coarsely segments garbage in geo-tagged photos so that citizens can report it. Kumar et al.~\cite{kumar2021novel} motivate their YOLOv3 and YOLOv3-tiny study with Indian waste statistics. They use 6,437 self-collected images in six classes and report 94.99\% and 51.95\% mAP on a split that served for both validation and testing, but the images are mostly clean, show one or a few items each, and their location is not stated. AquaVision~\cite{panwar2020aquavision} detects waste in water bodies using a subset derived from TACO, and TrashBox~\cite{kumsetty2022trashbox}, from NIT Karnataka, is a classification set.

The closest thing to an Indian street benchmark is DWSD~\cite{ali2025dwsd}: 784 phone images of streets, parks, lawns and garbage areas around Kolkata (Jadavpur, Garia) and Berhampore, West Bengal, with pixel masks for 14 waste classes. Its baselines report pixel accuracy only (89.39\% for DeepLabv3+) and no detection results, and the paper is not consistent about the dataset's name and image counts. Elsewhere in South Asia, UGV-NBWASTE~\cite{isalm2025ugvnbwaste} has 3,600 images of eight non-biodegradable waste types from Barisal, Bangladesh, with oriented boxes but only about 1.1 objects per image and no baselines. The pLitterStreet preprint~\cite{mandhati2024plitterstreet} maps street-level plastic litter from vehicle-mounted cameras in Thailand, Vietnam and Sri Lanka. Fang et al.~\cite{fang2023artificial} list the shortage of data in developing countries as one of the field's limits.

\subsection{Hyperparameter Optimization and Swarm Intelligence}
\label{subsec:rw-hpo}

Particle swarm optimization~\cite{kennedy1995particle} moves a population of candidate solutions, called particles, through the search space. The velocity of each particle combines its previous motion with a pull toward its own best position and a pull toward the best position of the swarm. Shi and Eberhart~\cite{shi1998modified} added the inertia weight, a factor on the previous velocity, and Clerc and Kennedy~\cite{clerc2002particle} derived coefficients that control convergence, which led to constriction-factor PSO. PSO needs no gradients and evaluates a fixed number of candidates per round, which suits expensive evaluations such as network trainings. Our settings are given in Section~\ref{subsec:pso}.

Random search~\cite{bergstra2012random} is the obvious reference for any guided search. It samples settings independently, so if an optimizer cannot beat it at the same number of evaluations, the optimizer's guidance has not been shown to help. Multi-fidelity methods take another route and make each evaluation cheaper. Hyperband~\cite{li2018hyperband}, for example, gives more resources to promising configurations and stops weak ones early. We use a simpler fixed low fidelity, a short proxy training (Section~\ref{subsec:proxy}).

Metaheuristic tuning of CNNs has been reported in several fields, but the evidence we found is thin. Aguerchi et al.~\cite{aguerchi2024cnn} use PSO to choose three architecture hyperparameters (filters, kernel size and stride) of a small CNN for mammogram masses, with the training hyperparameters fixed. Four particles run for three iterations, about 12 network evaluations. There is no baseline of any kind and one reported result, although the authors themselves note that a metaheuristic needs repeated runs and statistical analysis. The most complete comparison we found in the waste domain is by Sayed et al.~\cite{sayed2024intelligent}. A multi-objective beluga whale optimizer (MBWO) tunes the learning rate, batch size and dropout period of InceptionV3 on TrashNet, with a population of 30, 20 iterations and 15 epochs per evaluation. Four other metaheuristics are run with the same population, iterations, bounds and fitness, and an untuned model (90.76\% accuracy) serves as a reference. A Wilcoxon rank-sum test on the best search fitness over independent runs finds MBWO significantly better than three competitors but not the fourth (EO). The number of runs is not stated, random search is absent, the test metrics are single numbers without intervals, and the headline accuracy is 97.75\% in the abstract but 97.57\% in the results table. Mao et al.~\cite{mao2021recycling} apply a genetic algorithm to a CNN that classifies recycling waste.

\subsection{Comparison and Research Gaps}
\label{subsec:gaps}

Table~\ref{tab:lit-matrix} compares the 15 studies we read in full on the points that decide whether a tuning claim can be trusted. Only two of them tune hyperparameters with a metaheuristic, both for classifiers, and only one of those compares tuning methods. None reports a confidence interval on a test metric. Only Bartolo et al.~\cite{bartolo2025learning} test on public datasets not seen in training. From this table and the sections above we draw four gaps.

\begin{table*}[!t]
\centering
\caption{Comparison of the Studies Read in Full With This Work}
\label{tab:lit-matrix}
\footnotesize
\setlength{\tabcolsep}{3pt}
\renewcommand{\arraystretch}{1.1}
\begin{tabular}{@{}>{\raggedright\arraybackslash}p{2.75cm}>{\raggedright\arraybackslash}p{1.45cm}>{\raggedright\arraybackslash}p{1.75cm}>{\raggedright\arraybackslash}p{2.2cm}>{\raggedright\arraybackslash}p{2.75cm}>{\raggedright\arraybackslash}p{2.45cm}>{\raggedright\arraybackslash}p{1.95cm}>{\raggedright\arraybackslash}p{1.3cm}@{}}
\toprule
Study & Task & Scenes & Hyperparameters set by & Tuning baselines & Runs / CIs & Cross-domain test & India data \\
\midrule
Yang \& Thung~\cite{yang2016classification} & Cls & Studio & Manual & --- & 1 / no & No & No \\
Kaya~\cite{kaya2023classification} & Cls & TrashNet & Fixed & --- & 1 / no & No & No \\
Sayed et al.~\cite{sayed2024intelligent} & Cls & TrashNet & MBWO & 4 metaheur.\ (equal), untuned; no RS & n.r.\ + Wilcoxon (fitness) / no & No & No \\
Aguerchi et al.~\cite{aguerchi2024cnn} & Cls & Clinical & PSO (architecture) & None & 1 / no & No & --- \\
Majchrowska et al.~\cite{majchrowska2022deep} & Det + Cls & Real-world & Manual & None & 1 / no & No & No \\
Hong et al.~\cite{hong2020trashcan} & Det / Seg & Underwater & Fixed & None & 1 / no & No & No \\
Kraft et al.~\cite{kraft2021autonomous} & Det & Aerial & Not described & None & 1 / no & No & No \\
C\'{o}rdova et al.~\cite{cordova2022litter} & Det & Real-world & Published defaults & None & No SDs / no & Per dataset & No \\
Wang et al.~\cite{wang2020multilevel} & Seg & Lab + TACO & Fixed; CRF grid & None & 1 / no & No & No \\
Fulton et al.~\cite{fulton2019robotic} & Det & Underwater & Community & None & 1 / no & Shifted footage & No \\
Melinte et al.~\cite{melinte2020deep} & Det & TrashNet & Manual + LR test & None & 1 / no & Lab-made set & No \\
Kumar et al.~\cite{kumar2021novel} & Det & Mostly clean & Fixed & None & 1 / test = val & No & Unstated \\
Liao \& Juang~\cite{liao2022realtime} & Det & Aerial & ``Tuned'' (n.r.) & None & No & Qualitative & No \\
Bartolo et al.~\cite{bartolo2025learning} & Det & Aerial & Manual & None & 1 / no & Yes & No \\
Ali et al.~\cite{ali2025dwsd} & Seg & Real-world & n.r. & --- & 1 / no & No & Yes \\
\midrule
\textbf{This work} & \textbf{Det} & \textbf{Studio + real-world} & \textbf{PSO vs RS vs defaults} & \textbf{Equal: 32 proxy trainings each} & \textbf{1 seed / bootstrap 95\% CIs} & \textbf{Yes (3 domains)} & \textbf{DWSD (test)} \\
\bottomrule
\end{tabular}

\vspace{3pt}
\begin{minipage}{\textwidth}
\scriptsize
Datasets: TrashNet~\cite{yang2016classification} (single-object studio images), TrashCan~\cite{hong2020trashcan}, UAVVaste~\cite{kraft2021autonomous}, PlastOPol~\cite{cordova2022litter}, MJU-Waste~\cite{wang2020multilevel} and DWSD~\cite{ali2025dwsd} (a dataset paper). Cls = classification; Det = detection (one class in~\cite{majchrowska2022deep,cordova2022litter}); Seg = segmentation; MBWO = multi-objective beluga whale optimizer; metaheur.\ (equal) = other metaheuristics with the same population size and iterations; RS = random search; CRF grid = grid search for the conditional random field only; LR test = learning-rate range test; SD = standard deviation; CI = confidence interval; n.r.\ = not reported; --- = not applicable. ``1 / no'' = one reported run per configuration and no confidence intervals. ``Fixed'' = one setting used without a described selection procedure. ``test = val'' = the split used for validation also served as the test split. Per dataset = two datasets, each trained and tested separately; Qualitative = new sites described only qualitatively; Unstated = Indian authors, image location not stated. This work: 6 classes in \genone and 7 in \gentwo; one training seed per arm; paired-bootstrap CIs; studio, real-world and Indian test sets. Full-text works not listed are the review of Fang et al.~\cite{fang2023artificial} and the UGV-NBWASTE dataset~\cite{isalm2025ugvnbwaste}, which reports no model; works read only as abstracts or records are discussed in the text.
\end{minipage}
\end{table*}

\textit{Gap 1 (task and scene).} The classifiers of Section~\ref{subsec:rw-waste}~\cite{yang2016classification,kaya2023classification,sayed2024intelligent,alshammari2025smart,kumsetty2022trashbox} give one label per image, and the TrashNet-based ones see one object on a plain background. Fang et al.~\cite{fang2023artificial} find whole-image classification to be the most common approach, and even some detection studies use data with nearly one object per image~\cite{melinte2020deep,kumar2021novel,isalm2025ugvnbwaste}. Multi-object waste detection is still an established line of work, with the datasets and pipeline described above~\cite{proenca2020taco,bashkirova2022zerowaste,cheng2021flow,yudin2024hierarchical,kraft2021autonomous,cordova2022litter,hong2020trashcan,majchrowska2022deep}, and since our review is not systematic we do not estimate how much smaller it is. The gap, in the papers we reviewed, is that waste recognition for sorting is mostly single-label classification of one object on a plain background, that multi-object detection in cluttered scenes is the smaller part, and that several detection studies still depend on single-object studio images.

\textit{Gap 2 (optimization).} Metaheuristic tuning in waste vision has been applied to classifiers~\cite{sayed2024intelligent,mao2021recycling}, and PSO tuning of CNNs appears in other fields~\cite{aguerchi2024cnn}. None of the waste detection or segmentation studies we read in full describes a search over training hyperparameters (Table~\ref{tab:lit-matrix}). The only search that is described covers post-processing parameters~\cite{wang2020multilevel}, and in the sections of Fang et al.~\cite{fang2023artificial} that we read, PSO and genetic algorithms come up only for routing and prediction. Two caveats apply. Two detection papers say that they tuned their models without saying how~\cite{liao2022realtime,bartolo2025learning}, so a search there cannot be ruled out, and we could not check the training setup of papers known to us only from the abstract. With those caveats, the gap is that no waste detector in the papers we reviewed had its training hyperparameters tuned by a swarm optimizer. Detection studies used default or hand-set values, or a tuning step they do not describe.

\textit{Gap 3 (evidence).} Where tuning is evaluated, the baselines are weak: none in Aguerchi et al.~\cite{aguerchi2024cnn} and only other metaheuristics in Sayed et al.~\cite{sayed2024intelligent}, with single-number test metrics in both (Section~\ref{subsec:rw-hpo}). No detection or segmentation paper we read in full reports run-to-run variation, confidence intervals or significance tests on its accuracy (Table~\ref{tab:lit-matrix}), and Bartolo et al.~\cite{bartolo2025learning} state that statistical analysis was not applicable. Sayed et al.\ do compare optimizers at equal population and iterations, with an untuned reference and a significance test on the search fitness, so we do not claim that nobody compares optimizers or tests significance. What we did not find is a comparison with random search~\cite{bergstra2012random} at an equal evaluation budget, or test results reported with confidence intervals. Our own study narrows this gap but does not close it, because with one training seed per arm and one search seed per method our intervals cover test-set sampling only (Section~\ref{subsec:limitations}).

\textit{Gap 4 (generalization).} The TrashNet-based classifiers never leave the studio, two reviews point to simplified, lab-like data~\cite{lu2022computer,fang2023artificial}, and accuracy changes a lot from one dataset to another (the PlastOPol example in Section~\ref{sec:intro}~\cite{cordova2022litter}). Most detection studies test within one dataset or describe the drop at new sites only qualitatively (Table~\ref{tab:lit-matrix}), and the Indian work of Section~\ref{subsec:rw-india} offers no detection benchmark. Cross-domain tests are not entirely absent. Bartolo et al.~\cite{bartolo2025learning} test on datasets unseen in training, Fulton et al.~\cite{fulton2019robotic} on footage from a different time and place, and S\'{a}nchez-Ferrer et al.~\cite{sanchezferrer2023experimental} add synthetic scenes to scarce real data. South Asian detection data exist in UGV-NBWASTE~\cite{isalm2025ugvnbwaste}, without baselines, and in the pLitterStreet preprint~\cite{mandhati2024plitterstreet}. The gap is that studio-to-real-world and cross-dataset evaluation is rare and recent, and that we found no Indian real-world waste test set with published detection baselines. We report studio, real-world and DWSD results separately and compare studio-only with mixed training on the same test images (Section~\ref{subsec:res-domain}).

One more pattern is worth noting, although a targeted review cannot prove it. The applied papers we read present their methods as improvements. Sayed et al.~\cite{sayed2024intelligent}, for example, call their optimizer the best even though its edge over EO is not significant. Some negative findings are reported: a CNN far below a support vector machine~\cite{yang2016classification}, multi-class detection set aside because it worked poorly~\cite{majchrowska2022deep}, synthetic data that helped only a little~\cite{sanchezferrer2023experimental}, and weaker results on shifted data~\cite{fulton2019robotic,liao2022realtime}. We found no paper reporting that a tuning method failed to beat a simple baseline, which is what we report here. Table~\ref{tab:gap-map} lists how this study responds to each gap.

\begin{table}[!t]
\centering
\caption{Research Gaps and How This Study Addresses Them}
\label{tab:gap-map}
\footnotesize
\setlength{\tabcolsep}{3pt}
\renewcommand{\arraystretch}{1.1}
\begin{tabular}{@{}>{\raggedright\arraybackslash}p{1.45cm}>{\raggedright\arraybackslash}p{4.85cm}>{\raggedright\arraybackslash}p{1.95cm}@{}}
\toprule
Gap & How this study addresses it & Where \\
\midrule
G1: task and scene & Multi-object detection with YOLOv8n (6 or 7 classes) on studio and real-world scenes & Sec.~\ref{subsec:datasets}, \ref{subsec:detector} \\
G2: optimization & A hand-written PSO tunes 6 training hyperparameters of a waste detector & Sec.~\ref{subsec:space}, \ref{subsec:pso} \\
G3: evidence & Defaults, random search and PSO at 32 proxy trainings each; proxy check with a threshold fixed in advance (\genone only); test split used once; paired-bootstrap 95\% CIs & Sec.~\ref{subsec:design}, \ref{subsec:proxy}, \ref{subsec:eval}, \ref{subsec:res-gen1}, \ref{subsec:res-gen2} \\
G4: generalization & Real-world training data in \gentwo; studio, real-world and India (DWSD, test only, boxes from masks) results reported separately; both generations on the same test images & Sec.~\ref{subsec:datasets}, \ref{subsec:res-domain} \\
\bottomrule
\end{tabular}
\end{table}
